\chapter{Exact Reconstruction: Continuous RSVP $\to$ Polyxan Hypergraphs}
\label{ch:reconstruction}

\section*{Overview}

This chapter establishes the inverse direction of the semantic correspondence.
Where Chapter~28 constructed a canonical functor
\[
\mathcal{S} : \mathbf{Polyx} \longrightarrow \mathbf{RSVP},
\]
sending a finite typed Polyxan hypergraph to its stratified RSVP embedding,
we now construct the exact inverse
\[
\mathcal{R} : \mathbf{RSVP} \longrightarrow \mathbf{Polyx},
\]
and prove that $\mathcal{R}$ is a strict inverse on the nose:
\[
\mathcal{R}\circ\mathcal{S} \;\cong\; \mathrm{id}_{\mathbf{Polyx}},
\qquad
\mathcal{S}\circ\mathcal{R} \;\cong\; \mathrm{id}_{\mathbf{RSVP}}^{\mathrm{harm}},
\]
where the second isomorphism holds on the full subcategory of
harmonic RSVP configurations.

Thus the continuous and discrete semantic universes are
\emph{exactly equivalent}.  
Every RSVP field configuration \emph{is} a Polyxan hypergraph in disguise.

\section{Stratified RSVP Data}

Let $(\Phi,\mathbf{v},S)$ be a smooth RSVP field configuration on $\mathcal{M}$
satisfying:

\begin{enumerate}
\item $\mathcal{R} = \Delta\Phi - \beta\,\nabla\cdot\mathbf{v}$  
      is piecewise-constant on a finite stratification
      \[
      \mathcal{M} = \bigsqcup_{k=0}^N \mathcal{M}_k
      \]
      with locally-finite boundaries.

\item $S$ is constant on each $\mathcal{M}_k$.

\item $\iota : \bigsqcup_k \mathcal{M}_k \to \mathcal{M}$ is the identity map,
      but each stratum $\mathcal{M}_k$ is labeled by its curvature and entropy
      values.
\end{enumerate}

Such configurations form the full subcategory
$\mathbf{RSVP}^{\mathrm{strat}} \subset \mathbf{RSVP}$.

We show that every stratified configuration determines a unique
typed Polyxan hypergraph.

\section{Cells From Curvature Strata}

Define the $k$-cells of the reconstructed hypergraph by:
\[
c_k = \text{connected components of } \mathcal{M}_k.
\]

Each cell inherits two invariants:

\begin{enumerate}
\item \textbf{Type:}
      \[
      \mathsf{type}(c_k)
      := \mathcal{R}|_{c_k} \in \mathbb{Z}_{\ge 0}
      \]
      (curvature stratum index, matching Chapter~13).

\item \textbf{Tag:}
      \[
      \mathsf{tag}(c_k)
      := S|_{c_k} \in \mathbb{N}
      \]
      (entropy stratum index, matching Chapters~24–26).
\end{enumerate}

Incidence relations are induced directly from stratum boundaries:
\[
c_k \prec c_{k-1}
\quad\Longleftrightarrow\quad
\overline{c_k} \cap c_{k-1} \neq \varnothing.
\]

Thus the entire cell complex of $\mathcal{G} = \mathcal{R}(\Phi,\mathbf{v},S)$
is recovered from the stratified geometry.

\section{Hyperedges From Vector Field Flowlines}

The RSVP vector field $\mathbf{v}$ determines the hyperedges.

Define a \emph{hyperedge} as an equivalence class of integral curves of $\mathbf{v}$
flowing from one stratum to another:
\[
e = [\gamma] \quad\text{where}\quad \dot{\gamma} = \mathbf{v}(\gamma).
\]

Endpoints of $\gamma$ define the incident cells.

Types and tags of hyperedges are inherited from the strata traversed:
\[
\mathsf{type}(e) = \max_{t}\mathcal{R}(\gamma(t)),
\qquad
\mathsf{tag}(e) = \max_{t}S(\gamma(t)).
\]

In harmonic configurations, $\mathbf{v}$ foliation is globally trivial,
ensuring the finiteness of the hyperedge set.

\section{Nodes From Equilibrium Points}

Equilibrium points of the RSVP flow
\[
\nabla\Phi(x)=0,\qquad \mathbf{v}(x)=0,\qquad \nabla S(x)=0
\]
are declared as \emph{nodes}.

Each node inherits its type and tag from the stratum in which it sits.

These nodes correspond precisely to the atomic units $\mathcal{G}_{\mathrm{atomic}}$
used in the quine tower of Chapter~25.

\section{Reconstruction Functor}

We now define the reconstruction functor
\[
\mathcal{R} : \mathbf{RSVP}^{\mathrm{strat}} \longrightarrow \mathbf{Polyx}
\]
by:
\begin{enumerate}
\item Objects:  
      $\mathcal{R}(\Phi,\mathbf{v},S)$ is the hypergraph with:
      \[
      \mathrm{Cells} = \{c_k\},
      \qquad
      \mathrm{Edges} = \{e\},
      \qquad
      \mathrm{Nodes} = \{x\}.
      \]

\item Morphisms:  
      A stratified RSVP map $f:\mathcal{M}\to\mathcal{N}$ preserving
      $(\Phi,\mathbf{v},S)$ on strata induces a unique typed-hypergraph
      morphism preserving incidence, types, and tags.
\end{enumerate}

\section{Exactness on Harmonic Configurations}

\begin{theorem}[Exactness]
If $(\Phi,\mathbf{v},S)$ is harmonic, then:
\[
\mathcal{S}\circ\mathcal{R}(\Phi,\mathbf{v},S)
\;\cong\;
(\Phi,\mathbf{v},S)
\]
canonically.
\end{theorem}

\begin{proof}
In harmonic configurations:

\begin{enumerate}
\item Curvature strata are eigen-level surfaces of $\Phi$.  
\item Entropy strata are eigen-level surfaces of $S$.  
\item Flowlines of $\mathbf{v}$ foliate each stratum without tangential twist.  
\end{enumerate}

Thus the entire stratification is rigid, and the reconstruction followed by
semantic embedding reproduces exactly the original configuration.
\end{proof}

\section{Equivalence of Universes}

\begin{theorem}[Equivalence of Categories]
The functors $\mathcal{S}$ and $\mathcal{R}$ form an equivalence:
\[
\mathcal{S}:\mathbf{Polyx}\;\rightleftarrows\;\mathbf{RSVP}^{\mathrm{harm}}:\mathcal{R}.
\]
\end{theorem}

\begin{proof}
Ch.~28 shows $\mathcal{S}$ is fully faithful and essentially surjective.  
This chapter proves $\mathcal{R}$ is a strict inverse on harmonic objects.  
Thus $\mathcal{S}$ is an equivalence of categories with inverse $\mathcal{R}$.
\end{proof}

\section{Final Cadence}

\begin{quote}
\textbf{
Every continuous configuration is secretly discrete.  
Every geometric curvature pattern is secretly a quine.  
Every RSVP universe is a Polyxan hypergraph in disguise.
}
\end{quote}

The duality between the continuous and discrete semantic universes is exact.  
There is no residue, no loss, no approximation.

The universe of meaning is one.


